\section{Evaluation}
\label{sec:evaluation}
We compare the three representations on three patterns: typing at the end
of a file, typing in the middle of a large file, and replacing every
occurrence of a word. Table~\ref{tab:costs} summarizes the cost of a single
edit at distance $\delta$ from the previous one in a file of $n$ characters.

\begin{table}[h]
  \centering
  \begin{tabular}{lccc}
    & typing at the end & typing in the middle & reading \\
    \hline
    gap buffer  & $O(1)$       & $O(\delta)$ & $O(1)$ \\
    piece table & $O(1)$       & $O(p)$      & $O(p)$ \\
    rope        & $O(\log n)$ & $O(\log n)$ & $O(\log n)$ \\
  \end{tabular}
  \caption{Cost of one edit; $p$ is the number of pieces.}
  \label{tab:costs}
\end{table}

When typing at the end of a file, all three representations perform well.
The gap buffer is fastest, because the gap is already at the cursor and
each key press writes one byte.

When typing in the middle of a large file, the gap buffer must first move
the gap, which costs time proportional to the distance from the previous
cursor position. The piece table and the rope are unaffected by the
distance.

Replacing every occurrence of a word makes many edits spread across the
document. The rope handles each edit in logarithmic time and keeps every
intermediate version, while the piece table accumulates many small pieces
that slow down later reads until it is compacted.

Taken together, these patterns suggest that ropes are a good default for
editors that open large files, and that gap buffers remain attractive for
small files and simple editors.
